% modD_control.tex — module D: predictive test-time control.
% Draws eq:rollout (latent rollout per candidate, no LLM call), the
% closed-form expansion's per-step command injection (eq:closed_form_rollout),
% and eq:selection (cost + argmin over the candidate set).
% Canvas 4.735 x 3.03; top-left 2.2 x 0.30 left empty for the corner tag.
\path (0,0) rectangle (4.735,3.03);

% ---------------- 1. candidate set C_t --------------------------------
\node[mathlab,anchor=south] at (0.51,2.57) {$\symCmdSet$};
\ScalarCell{(0.12,2.55)}{ours}{}
\ScalarCell{(0.40,2.55)}{bad}{}
\ScalarCell{(0.68,2.55)}{good}{}

% three thin lines: the same chain is rolled out once per candidate
\draw[draw=ours,line width=0.4pt] (0.23,2.33) -- (0.14,2.20);
\draw[draw=bad,line width=0.4pt]  (0.51,2.33) -- (0.14,2.20);
\draw[draw=good,line width=0.4pt] (0.79,2.33) -- (0.14,2.20);

% ---------------- 2. latent rollout without an LLM call ---------------
% xihat_{t+h|t} = K xihat_{t+h-1|t} + B c + b   (eq:rollout, eq:closed_form_rollout)
% label the chain ONCE at each end; the fading colour marks successive h.
\LatCol{(0.05,2.20)}{3}{0.18}
\node[mathlab,anchor=north] at (0.14,1.62) {$\symLat$};
\LatColVF{(0.69,2.20)}{3}{0.18}{78}
\LatColVF{(1.33,2.20)}{3}{0.18}{58}
\LatColVF{(1.97,2.20)}{3}{0.18}{42}
\node[mathlab,anchor=north] at (2.06,1.62) {$\symLatRoll$};

\draw[thinflow] (0.25,1.93) -- (0.67,1.93);
\draw[thinflow] (0.89,1.93) -- (1.31,1.93);
\node[mathlab,anchor=south] at (1.10,1.99) {$\symK$};
\draw[thinflow] (1.53,1.93) -- (1.95,1.93);

% one instance of the per-step command injection is enough: the three
% arrows are identical, so a single Bc (ours) + b (grey) reads as "every
% step adds the command".
\MatBlock{(1.02,1.86)}{3}{1}{0.07}{ours}
\MatBlock{(1.10,1.86)}{3}{1}{0.07}{black}

\DimBraceB{(0.05,1.30)}{(2.15,1.30)}{$\symStep = 1,\ldots,\symHor$}
\node[minicap] at (1.10,0.85) {no LLM call};

% ---------------- 3. decode to an attribute (readout, once) -----------
% sits on the chain's own line, to the right of it -- not stacked below.
\draw[thinflow] (2.17,1.93) -- (2.36,1.83);
\begin{scope}[shift={(2.42,1.80)},scale=0.26]
  \DecNet{mdDec}{(0,0)}
\end{scope}
\draw[thinflow] (2.72,1.80) -- (2.75,1.78);
\RowSelector{(2.75,1.87)}{3}{0.06}
\node[mathlab,anchor=south] at (2.84,1.87) {$\symReadout$};
\draw[thinflow] (2.93,1.78) -- (2.95,1.55);

% ---------------- 4. predicted trajectory vs. target -------------------
\MiniAxes{(2.95,0.55)}{1.50}{1.10}
\node[mathlab,anchor=west] at (4.47,0.55) {$\symStep$};
\node[mathlab,anchor=south east] at (2.98,1.66) {$\symYroll$};
\draw[dashed,draw=good,line width=0.5pt] (2.95,1.10) -- (4.45,1.10);
\node[mathlab,anchor=west,text=good] at (4.47,1.10) {$\symTarget$};
\draw[draw=bad,line width=0.6pt]  (3.00,0.70) -- (3.45,1.40) -- (3.90,1.50) -- (4.40,1.35);
\draw[draw=good,line width=0.6pt] (3.00,0.70) -- (3.45,0.95) -- (3.90,1.05) -- (4.40,1.10);
\draw[draw=ours,line width=0.6pt] (3.00,0.70) -- (3.45,0.78) -- (3.90,0.85) -- (4.40,0.92);
\node[minicap] at (4.47,0.92) {$\symBound$};

% ---------------- 5. cost and the choice (above the plot) --------------
\node[mathlab,anchor=south] at (3.30,2.66) {$\symCost$};
\fill[ours] (2.95,2.48) rectangle (3.50,2.60);
\fill[bad]  (2.95,2.30) rectangle (3.80,2.42);
\fill[good] (2.95,2.12) rectangle (3.20,2.24);
\ScalarCell{(4.15,2.24)}{ours}{$\symCmdStar$}
\draw[oursflow] (4.15,2.13) -- (3.20,2.18);

% ---------------- 6. receding horizon (routed clear of the plot/bars) --
\draw[thinflow] (4.26,2.02) -- (4.70,2.02) -- (4.70,0.15) -- (0.06,0.15) -- (0.06,0.33);
\node[minicap,anchor=south] at (2.30,0.16) {execute, observe, replan};
